% Główny plik dokumentu – kompiluj właśnie ten plik

\documentclass[a4paper,12pt,oneside,onecolumn]{book}
% Ustawienia formatu dokumentu:
% - a4paper: format A4
% - 12pt: wielkość czcionki
% - oneside: jednostronny druk
% - onecolumn: jedna kolumna tekstu

\linespread{1.5} % Zwiększenie interlinii (1.5x) dla lepszej czytelności
% Kodowanie i czcionka
\usepackage[utf8]{inputenc}         % Kodowanie UTF-8
\usepackage[T1]{fontenc}            % Lepsze dzielenie wyrazów i znaki diakrytyczne
\usepackage{helvet}                 % Czcionka Helvetica
\renewcommand{\familydefault}{\sfdefault} % Domyślna czcionka bezszeryfowa

% Matematyka i symbole
\usepackage{amsmath,amsfonts,amssymb,amsthm} % Rozszerzenia matematyczne


% Marginesy
%\usepackage[margin=2.5cm]{geometry} % Ustawienie marginesów
\usepackage[top=2.5cm, bottom=2.5cm, left=3.5cm, right=1.5cm]{geometry} % lewy i prawwy inaczej dla marginesow 


% Ulepszenia typografii
\usepackage{microtype}
\usepackage{xcolor}
\usepackage{array}
\usepackage{enumitem} % opisy do wzorow
% ODSTEPY ENUMERATE 1. 2. 3. 
\setlist[enumerate]{itemsep=0pt, parsep=0pt, topsep=0pt}

% Odnośniki i linki
\usepackage{url}
\usepackage[hidelinks]{hyperref}   % Linki bez kolorów/ramki
\urlstyle{same} % linki taka sama czcionka



% --- FORMATOWANIE NAGŁÓWKÓW (Musi być dodane) ---
\usepackage{titlesec}

% Formatowanie ROZDZIAŁÓW (Nagłówek 1)
% Wymogi: 16 pkt, pogrubiona, odstępy 18pkt przed i po
\titleformat{\chapter}[block]
  {\bfseries\fontsize{16}{16}\selectfont} % Styl: Pogrubienie, rozmiar 16pt
  {\chaptername\ \thechapter.} % Numeracja: "Rozdział X."
  {0.5em} % Odstęp między "Rozdział X." a tytułem
  {}

% Odstępy dla rozdziału: {lewy}{góra}{dół}
\titlespacing*{\chapter}{0pt}{-20pt}{18pt}

% Formatowanie ROZDZIAŁÓW nienumerowanych (Spisy, Literatura, Wstęp)
\titleformat{name=\chapter,numberless}[block]
  {\bfseries\fontsize{16}{16}\selectfont}
  {}
  {0pt}
  {}
\titlespacing*{name=\chapter,numberless}{0pt}{-20pt}{18pt}

% Formatowanie PODROZDZIAŁÓW (Nagłówek 2 - np. 1.1)
% Wymogi: 14 pkt, pogrubiona, odstępy 12pkt przed i po
\titleformat{\section}[block]
  {\bfseries\fontsize{14}{14}\selectfont}
  {\thesection.}
  {0.5em}
  {}
\titlespacing*{\section}{0pt}{12pt}{12pt}

% Formatowanie PODPUNKTÓW (Nagłówek 3 - np. 1.1.1)
% Wymogi: 12 pkt, pogrubiona, odstępy 6pkt przed i po
\titleformat{\subsection}[block]
  {\bfseries\fontsize{12}{12}\selectfont}
  {\thesubsection.}
  {0.5em}
  {}
\titlespacing*{\subsection}{0pt}{6pt}{6pt}



% Wcięcia i bibliografia ************************************************************************************
\usepackage{indentfirst}           % Wcięcie w pierwszym akapicie każdego rozdziału
\setlength{\parindent}{1cm}

% --- BIBLIOGRAFIA I CYTOWANIA (Poprawiona) ---
\usepackage[polish]{babel}
\usepackage{csquotes}
\MakeOuterQuote{"}

% Konfiguracja BibLaTeX - styl authoryear (APA-like)
\usepackage[
    style=authoryear,    % Styl: Autor (Rok)
    backend=biber,       % Silnik przetwarzania
    maxcitenames=2,      % W tekście: max 2 nazwiska, potem "i in." (zgodnie z APA)
    maxbibnames=99,      % W bibliografii: wymień wszystkich autorów
    uniquelist=false,    % Nie dodawaj inicjałów, jeśli nazwiska się powtarzają
    isbn=false,          % WYŁĄCZ ISBN (zgodnie z szablonem)
    doi=false,           % WYŁĄCZ DOI (zgodnie z szablonem)
    url=false,           % WYŁĄCZ URL dla książek/artykułów
    eprint=false
]{biblatex}

% Włączamy URL tylko dla typu @online (netografia)
\AtEveryBibitem{%
  \ifentrytype{online}{}{\clearfield{url}}%
  \ifentrytype{online}{}{\clearfield{urldate}}%
}

% --- FORMATOWANIE WYGLĄDU bibliografi ---

% 1. Wszystkie tytuły kursywą (książki, artykuły, strony www)
\DeclareFieldFormat{title}{\mkbibemph{#1}}          
\DeclareFieldFormat[article]{title}{\mkbibemph{#1}} 
\DeclareFieldFormat[book]{title}{\mkbibemph{#1}}    
\DeclareFieldFormat[online]{title}{\mkbibemph{#1}}    

% 2. USUWANIE "W:" przed nazwą czasopisma
% Standardowo biblatex daje "In:" lub "W:". Ta komenda to usuwa.
\renewbibmacro{in:}{}

% 3. Kolejność: Wydawnictwo, Miasto (Dla książek)
% Standard to Miasto: Wydawnictwo. UE Katowice chce odwrotnie.
\renewbibmacro*{publisher+location+date}{%
  \printlist{publisher}%
  \setunit*{\addcomma\space}%
  \printlist{location}%
  \setunit*{\addcomma\space}%
  \usebibmacro{date}%
  \newunit}

% 4. Separatory
\renewcommand*{\nameyeardelim}{\addcomma\space} % Przecinek po nazwisku w cytowaniu
\renewcommand*{\postnotedelim}{\addcomma\space} % Przecinek przed numerem strony

% 5. Data dostępu dla netografii: [dostęp DD.MM.RRRR]
\DefineBibliographyStrings{polish}{
  urlseen = {dostęp},
}
\DeclareFieldFormat{urldate}{\mkbibbrackets{\bibstring{urlseen}\space#1}}
\DeclareFieldFormat{url}{\url{#1}}

% 6. LISTA NUMEROWANA (1. 2. 3.) do bibliografi 
\newcounter{mybibcounter}
\defbibenvironment{bibliography}
  {\list
     {\stepcounter{mybibcounter}\themybibcounter.} 
     {\setlength{\leftmargin}{2.5em}%              
      \setlength{\labelwidth}{2em}%                
      \setlength{\labelsep}{0.5em}%                
      \setlength{\itemindent}{0pt}%
      \setlength{\itemsep}{\bibitemsep}%
      \setlength{\parsep}{\bibparsep}}}
  {\endlist}
  {\item}





% Kod źródłowy (np. programy)
\usepackage{listings}

% Nagłówki i stopki
\usepackage{fancyhdr}
\usepackage{etoolbox}

% Grafika
\usepackage{graphicx}
\graphicspath{ {images/} }         % Folder z grafikami

% Opisy pod rysunkami
\usepackage{caption}
\usepackage{chngcntr}

% Globalne numerowanie rysunków i tabel
\counterwithout{figure}{chapter}
\counterwithout{table}{chapter}

%definicja podpisu obrazka
\DeclareCaptionFont{elevenpt}{\fontsize{11}{13.6}\selectfont} % Definicja 11pkt
\captionsetup{
  font=elevenpt, % Użycie 11pkt
  labelfont=bf,
  labelsep=period,
  justification=raggedright, 
  singlelinecheck=false
}


\usepackage{float}

\newcommand{\MyFigure}[5][0.7\textwidth]{%
  \begin{figure}[H] % H dla tego miejsca ht dla ustawienia roznie 
    \centering
    \includegraphics[width=#1]{#2}
    \caption{#3}
    \label{#4}
    \vspace{-3pt}
    {\raggedright\fontsize{10pt}{12pt}\selectfont\textit{Źródło: #5}\par}
  \end{figure}
}

% Niestandardowe środowisko do tabel z podpisem i źródłem ---- TABELE
\captionsetup[table]{
    format=plain,
    justification=raggedright,
   labelfont=bf,
   textfont=normalfont,
   labelsep=period,
   name=Tabela,
   singlelinecheck=false,
   skip=2pt
}

 \newcommand{\MyTableSource}[1]{%
   \vspace{-8pt}
   \begin{flushleft}
     {\linespread{0.9}\fontsize{10}{12pt}\selectfont\parbox{0.95\textwidth}{\raggedright\textit{Źródło: #1}}}
   \end{flushleft}
   \vspace{-6pt}
 }

% Nadpisanie stylu rozdziału – aby używał fancyhdr
\patchcmd{\chapter}{\thispagestyle{plain}}{\thispagestyle{fancy}}{}{}

% Stylowanie kodu źródłowego
\definecolor{codegreen}{rgb}{0,0.6,0}
\definecolor{codegray}{rgb}{0.5,0.5,0.5}
\definecolor{codepurple}{rgb}{0.58,0,0.82}
\definecolor{backcolour}{rgb}{0.98,0.98,0.95} % Nieco jaśniejsze tło

\lstdefinestyle{mystyle}{
    backgroundcolor=\color{backcolour},
    commentstyle=\color{codegreen},
    keywordstyle=\color{blue}\bfseries, % Słowa kluczowe na niebiesko i pogrubione
    numberstyle=\tiny\color{codegray},
    stringstyle=\color{codepurple},
    basicstyle=\ttfamily\footnotesize,
    breakatwhitespace=false,
    breaklines=true,
    captionpos=b,
    keepspaces=true,
    numbers=left,
    numbersep=10pt,                  % Większy odstęp numerów od kodu
    showspaces=false,
    showstringspaces=false,
    showtabs=false,
    tabsize=4,                       % Standard dla Pythona to 4 spacje
    frame=single,                    % Ramka wokół kodu
    rulecolor=\color{black!10},      % Kolor ramki
    extendedchars=true,              % Obsługa znaków specjalnych
    literate={ą}{{\k{a}}}1 {ć}{{\'c}}1 {ę}{{\k{e}}}1 {ł}{{\l{}}}1 {ń}{{\'n}}1 {ó}{{\'o}}1 {ś}{{\'s}}1 {ź}{{\'z}}1 {ż}{{\.z}}1 {Ą}{{\k{A}}}1 {Ć}{{\'C}}1 {Ę}{{\k{E}}}1 {Ł}{{\L{}}}1 {Ń}{{\'N}}1 {Ó}{{\'O}}1 {Ś}{{\'S}}1 {Ź}{{\'Z}}1 {Ż}{{\.Z}}1 % Polskie znaki w komentarzach
}

% Definicja konkretnie pod Pythona
\lstdefinestyle{pythonstyle}{
    style=mystyle,
    language=Python,
    morekeywords={self, cls, @property, @staticmethod, yield}
}

\lstset{style=pythonstyle} % Ustawienie jako domyślny

\lstdefinelanguage{javaScript}{
  keywords={typeof, new, true, false, catch, function, return, null, catch, switch, var, if, in, while, do, else, case, break},
  keywordstyle=\color{blue}\bfseries,
  ndkeywords={class, export, boolean, throw, implements, import, this},
  ndkeywordstyle=\color{darkgray}\bfseries,
  identifierstyle=\color{black},
  sensitive=false,
  comment=[l]{//},
  morecomment=[s]{/*}{*/},
  commentstyle=\color{purple}\ttfamily,
  stringstyle=\color{red}\ttfamily,
  morestring=[b]',
  morestring=[b]"
}

\lstdefinelanguage{css}{
  morekeywords={accelerator,azimuth,background,background-attachment,
    background-color,background-image,background-position,
    background-position-x,background-position-y,background-repeat,
    behavior,border,border-bottom,border-bottom-color,
    border-bottom-style,border-bottom-width,border-collapse,
    border-color,border-left,border-left-color,border-left-style,
    border-left-width,border-right,border-right-color,
    border-right-style,border-right-width,border-spacing,
    border-style,border-top,border-top-color,border-top-style,
    border-top-width,border-width,bottom,caption-side,clear,
    clip,color,content,counter-increment,counter-reset,cue,
    cue-after,cue-before,cursor,direction,display,elevation,
    empty-cells,filter,float,font,font-family,font-size,
    font-size-adjust,font-stretch,font-style,font-variant,
    font-weight,height,ime-mode,include-source,
    layer-background-color,layer-background-image,layout-flow,
    layout-grid,layout-grid-char,layout-grid-char-spacing,
    layout-grid-line,layout-grid-mode,layout-grid-type,left,
    letter-spacing,line-break,line-height,list-style,
    list-style-image,list-style-position,list-style-type,margin,
    margin-bottom,margin-left,margin-right,margin-top,
    marker-offset,marks,max-height,max-width,min-height,
    min-width,-moz-binding,-moz-border-radius,
    -moz-border-radius-topleft,-moz-border-radius-topright,
    -moz-border-radius-bottomright,-moz-border-radius-bottomleft,
    -moz-border-top-colors,-moz-border-right-colors,
    -moz-border-bottom-colors,-moz-border-left-colors,-moz-opacity,
    -moz-outline,-moz-outline-color,-moz-outline-style,
    -moz-outline-width,-moz-user-focus,-moz-user-input,
    -moz-user-modify,-moz-user-select,orphans,outline,
    outline-color,outline-style,outline-width,overflow,
    overflow-X,overflow-Y,padding,padding-bottom,padding-left,
    padding-right,padding-top,page,page-break-after,
    page-break-before,page-break-inside,pause,pause-after,
    pause-before,pitch,pitch-range,play-during,position,quotes,
    -replace,richness,right,ruby-align,ruby-overhang,
    ruby-position,-set-link-source,size,speak,speak-header,
    speak-numeral,speak-punctuation,speech-rate,stress,
    scrollbar-arrow-color,scrollbar-base-color,
    scrollbar-dark-shadow-color,scrollbar-face-color,
    scrollbar-highlight-color,scrollbar-shadow-color,
    scrollbar-3d-light-color,scrollbar-track-color,table-layout,
    text-align,text-align-last,text-decoration,text-indent,
    text-justify,text-overflow,text-shadow,text-transform,
    text-autospace,text-kashida-space,text-underline-position,top,
    unicode-bidi,-use-link-source,vertical-align,visibility,
    voice-family,volume,white-space,widows,width,word-break,
    word-spacing,word-wrap,writing-mode,z-index,zoom},
  morestring=[s]{:}{;},
  sensitive,
  morecomment=[s]{/*}{*/}
}

% Spolszczenia nazw elementów
\renewcommand{\lstlistingname}{Listing}
\renewcommand{\lstlistlistingname}{Spis listingów}
\renewcommand{\listfigurename}{Spis rysunków}
\renewcommand{\listtablename}{Spis tabel}
\renewcommand{\figurename}{Rysunek}
\renewcommand{\tablename}{Tabela}
\renewcommand{\chaptername}{Rozdział}
\renewcommand{\contentsname}{Spis Treści}

% Środowisko do warunków (np. w matematyce)
\newenvironment{conditions}
  {\par\vspace{\abovedisplayskip}\noindent\begin{tabular}{>{$}l<{$} @{${}={}$} l}}
  {\end{tabular}\par\vspace{\belowdisplayskip}}



% --- NOWE ŚRODOWISKO "GDZIE" (opis do wzorow) ---
\newenvironment{gdzie}
{
    \vspace{-1.0em}    % Podciągnij tekst do góry (bliżej wzoru)
    \noindent \textbf{gdzie:}  % Automatyczny napis "gdzie:"
    \begin{itemize}[         % Lista z ciasnymi ustawieniami
        noitemsep,           % Brak odstępów między punktami
        topsep=2pt,  % Mały odstęp od "gdzie:" do pierwszego punktu
        parsep=0pt,          % Brak odstępów wewnątrz punktów
        partopsep=0pt,       % Brak dodatkowych odstępów
        leftmargin=1.5em,    % Wcięcie listy
        labelsep=0.5em     % Odstęp kropki/myślnika od tekstu
    ]
}
{
    \end{itemize}
    \vspace{0.5em}                  % Odstęp po zakończeniu opisu
} % Dołączenie pliku z definicją stylu dokumentu

% Dołączenie plików z bibliografią (literatura i netografia)
\addbibresource{refs.bib}
\addbibresource{url_refs.bib}


\usepackage[titles]{tocloft}


\begin{document}

% *************** Front matter ***************
% *************** Strona tytułowa ***************

\pagestyle{empty}
\noindent
\begin{center}
    UNIWERSYTET EKONOMICZNY W KATOWICACH
\end{center}

\begin{center}
    INFORMATYKA
\end{center}

\vfill
\begin{center}
    \textbf{Patryk Klimek,}
    \textbf{Dominik Rekowski,}
    \textbf{Łukasz Osadnik}
\end{center}
\vfill

\begin{center}
    \Large
	\textbf{Aplikacja mobilna do kodowania "CodeMobile"}
\end{center}

\begin{center}
	\textbf{Mobile app for coding "CodeMobile"}
\end{center}

\vfill
\vfill
\vfill


\vfill
\begin{center}
    KATOWICE 2026
\end{center}

\newpage 

% *************** Spis treści ***************
\tableofcontents

% *************** Koniec front matter ***************

 % Tutaj znajduje się strona tytułowa, spis treści, oświadczenia itd.

% *************** Treść rozdziałów *******-> 80-130 tys znakow

% \chapter*{Wstęp}

\addcontentsline{toc}{chapter}{Wstęp}
\chapter{Aplikacje do programowania na telefonie}

Programowanie to jedna z najpopularniejszych aktywności informatyków naszego pokolenia na telefonie. Każda osoba w naszej grupie przynajmniej raz stała w korku wracając autobusem (na przykład z nie etycznej lekcji etyki) i zastanawiała się co by było gdyby mogła w tym momencie wyjąć telefon i zacząć programować. Innym przypadkiem kiedy aplikacja była by użyteczna jest moment kiedy chcemy stworzyć prosty prototyp strony a jednocześnie nie chcemy wyjść z łóżka, zamiast zabierać komputer do łóżka jak psychopata wygodniej było by sięgnąć po telefon i na szybko coś napisać. Są też ludzie (głównie dzieci) którzy nie mają dostępu do komputera z różnych powodów a i tak chcieli by coś tworzyć.
\newline\newline
Czy naprawdę powinniśmy tym biednym ludziom odbierać możliwość programowania. Naszym zdaniem nie, szczególnie w obecnej erze gdzie każdy może sobie wygenerować szczochy wykorzystując sztuczną inteligencje, powinniśmy wspierać ludzi którzy wciąż mimo przeciwności losu chcą stworzyć coś własnego zamiast iść na łatwiznę. Narządzie dla ludzi z prawdziwą pasją powinny być proste oraz użyteczne a naszą misją jest stworzenie właśnie takiej aplikacji.
% -*- TeX-master: "../main.tex" -*-
\chapter{Analiza konkurencji}

Przed rozpoczęciem opracowywania projektu naszej aplikacji, szczegółowo przeanalizowaliśmy rynek. Pobraliśmy i przetestowaliśmy programy konkurencji, oraz przeanalizowaliśmy opinie ich użytkowników. Zebrane informacje skompilowane są poniżej.

\section{Code App}
\MyFigure[0.3\textwidth]{images/konkurencja/codeapp.png}{Aplikacja Code App}{fig:codeapp}{App Store}
\begin{itemize}
    \item \textbf{Docelowa platforma:} iOS/IPadOS
    \item \textbf{Docelowa grupa:} Programiści 
    \item \textbf{Główne przewagi konkurencyjne:} Silnik Monaco (z VS Code). Posiada wbudowane lokalne środowiska (Node.js, Python, PHP, C++) oraz klienta Git. 
    \item \textbf{Słabe strony:} Ograniczenia wynikające z Sandbox (system blokuje uruchamianie zewnętrznych procesów, systemowych usług w tle i narzędzi do wirtualizacji), brak autouzupełniania nazw klas i importów dla C++ i Java, brak dokumentacji. 
    \item \textbf{Model biznesowy:} Jednorazowa opłata w wysokości 30 zł
\end{itemize}

\newpage
\section{AndroidIDE}
\MyFigure[0.3\textwidth]{images/konkurencja/androidIDE.png}{Aplikacja AndroidIDE}{fig:androidide}{Google Play}
\begin{itemize}
    
    \item \textbf{Docelowa platforma:} Android
    \item \textbf{Docelowa grupa:} Początkujący i zaawansowani programiści aplikacji na Androida 
    \item \textbf{Główne przewagi konkurencyjne:} Pełne wsparcie dla systemu Gradle, aplikacja open source, autouzupełnianie kodu Java/Kotlin, kompilacja, budowanie i uruchamianie aplikacji w pełni lokalnie, zintegrowany projektant UI. 
    \item \textbf{Słabe strony:} Wysokie wymagania sprzętowe, projekt został porzucony w 2024. 
    \item \textbf{Model biznesowy:} Projekt open source (Bezpłatna) 
\end{itemize}

\newpage
\section{Acode}
\MyFigure[0.3\textwidth]{images/konkurencja/acode.png}{Aplikacja Acode}{fig:acode}{Google Play}
\begin{itemize}
    
    \item \textbf{Docelowa platforma:} Android 
    \item \textbf{Docelowa grupa:} Web-deweloperzy (Front-end), programiści 
    \item \textbf{Główne przewagi konkurencyjne:} Możliwość otwierania dużych plików, autouzupełnianie i kolorowanie składni dla ponad 100 języków programowania, wbudowany podgląd stron internetowych, interaktywna konsola JavaScript, wiele open source wtyczek do pobrania w aplikacji, skróty klawiszowe. 
    \item \textbf{Słabe strony:} Mało intuicyjny design (ciężko się odnaleźć w aplikacji)
    \item \textbf{Model biznesowy:} Dostępna jest wersja bezpłatna z reklamami oraz wersja płatna bez reklam za 13,99 zł. 
\end{itemize}

\newpage
\section{Textastic}
\MyFigure[0.3\textwidth]{images/konkurencja/textastic.png}{Aplikacja Textastic}{fig:textastic}{App Store}
\begin{itemize}
    \item \textbf{Docelowa platforma:} iOS/IPadOS 
    \item \textbf{Docelowa grupa:} Web-developerzy 
    \item \textbf{Główne przewagi konkurencyjne:} Świetna obsługa FTP, SFTP, WebDAV, integracja z chmurą (Dropbox, Google Drive), wbudowany terminal SSH, możliwość podglądu stron internetowych, kolorowanie składni dla ponad 80 języków, uzupełnianie kodu dla HTML, CSS, JS, C i PHP.
    \item \textbf{Słabe strony:} Brak kompilatorów, brak integracji z Git, ograniczone autouzupełnianie (uzupełnianie w oparciu o słowa, a nie o semantykę jak np. w VSC) 
    \item \textbf{Model biznesowy:} Płatna: jednorazowy zakup za 399,99 zł lub comiesięczna subskrypcja za 14,99 zł. 
\end{itemize}

\newpage
\section{Spck Editor}
\MyFigure[0.3\textwidth]{images/konkurencja/spck.png}{Aplikacja Spck Editor}{fig:spck}{Google Play}
\begin{itemize}
    
    \item \textbf{Docelowa platforma:} Android / iOS
    \item \textbf{Docelowa grupa:} Front-end dev (HTML, CSS, JS) 
    \item \textbf{Główne przewagi konkurencyjne:} Integracja z Git, asystent AI, Silnik Monaco Editor, wbudowana przeglądarka z narzędziami deweloperskimi, Spck Labs oferujące udostępnione projekty innych użytkowników do nauki i zabawy. 
    \item \textbf{Słabe strony:} Asystent AI oraz Predictive keyboard dostępne tylko dla użytkowników z miesięczną subskrypcją. 
    \item \textbf{Model biznesowy:} Bezpłatna, oraz wersje płatne: Spck Lite w cenie 35,99 zł jednorazowo (dostajemy w niej ekskluzywny motyw, dostęp do Predictive keyboard, customowe Snippety i paleta poleceń), Spck for NodeJS jako subskrypcja za 4,49 zł/miesiąc (integracja z NodeJS), subskrypcja „Supporter” za 8,79 zł/miesiąc (całkowite usunięcie reklam, dostęp do wszystkich płatnych motywów wizualnych interfejsu oraz dodatkowych opcji personalizacji środowiska pracy) oraz subskrypcja „Gold” za 23,99 zł/miesiąc (zawiera wszystkie korzyści z wersji Supporter, a dodatkowo zapewnia stały przydział waluty „Gold” przeznaczonej na funkcje premium, takie jak inteligentne autouzupełnianie kodu wspierane przez sztuczną inteligencję oraz rozszerzone możliwości synchronizacji projektów). 
\end{itemize}

\newpage
\section{Pydroid 3} 
\MyFigure[0.3\textwidth]{images/konkurencja/pydroid.png}{Aplikacja Pydroid 3}{fig:pydroid}{Google Play}
\begin{itemize}
    
    \item \textbf{Docelowa platforma:} Android
    \item \textbf{Docelowa grupa:} Osoby uczące się języka Python, Entuzjaści Data Science i AI. 
    \item \textbf{Główne przewagi konkurencyjne:} W pełni gotowy do użytku po zainstalowaniu aplikacji, posiada własne repozytorium prekompilowanych bibliotek (pozwala na bezproblemową instalację ciężkich paczek takich jak NumPy, Pandas, SciPy, Matplotlib, a także na uruchomienie interaktywnych notatników Jupyter), Obsługa interfejsów graficznych (np. Tkinter), wbudowany kompilator C/C++ (dla zaawansowanych bibliotek Pythona) 
    \item \textbf{Słabe strony:} Agresywne reklamy, zaawansowane biblioteki (np. TensorFlow) wymaga wersji Premium aplikacji, uzupełnianie i sprawdzanie kodu tylko w wersji Premium. 
    \item \textbf{Model biznesowy:} Bezpłatna, oraz wersja Premium z podziałem na subskrypcję 29,99 zł/miesiąc lub jednorazowy zakup za 99,99 zł. 
\end{itemize}

\newpage
\section{Pyto}
\MyFigure[0.3\textwidth]{images/konkurencja/pyto.png}{Aplikacja Pyto}{fig:pyto}{App Store}
\begin{itemize}
    
    \item \textbf{Docelowa platforma:} iOS 
    \item \textbf{Docelowa grupa:} Programiści Pythona w ekosystemie Apple 
    \item \textbf{Główne przewagi konkurencyjne:} Pozwala na uruchamianie skryptów prosto z aplikacji Skróty, pozwala na kodowanie własnych widżetów na ekran główny, dostęp do ciężkich bibliotek (NumPy, Pandas, SciPy, Matplotlib itd.) od razu po pobraniu aplikacji, pozwala na tworzenie interfejsów systemu iOS. 
    \item \textbf{Słabe strony:} Niestabilność aplikacji, brak darmowej wersji, ubogi interfejs UI. 
    \item \textbf{Model biznesowy:} Wersja Lite za 29,99 zł oraz Wersja Pełna za 59,99 zł. 
\end{itemize}

\newpage
\section{Carnets - Jupyter}
\MyFigure[0.3\textwidth]{images/konkurencja/cartens.png}{Aplikacja Carnets - Jupyter}{fig:carnet}{App Store}
\begin{itemize}
    
    \item \textbf{Docelowa platforma:} iOS/IPadOS
    \item \textbf{Docelowa grupa:} Analitycy danych i badacze 
    \item \textbf{Główne przewagi konkurencyjne:} Jupyter Notebook w pełni lokalnie, potężne biblioteki analityczne (NumPy, Pandas, Matplotlib, Bokeh, Seaborn) od razu zainstalowane, możliwość edycji i zapisu notatników w chmurze. 
    \item \textbf{Słabe strony:} Potrzeba zainstalowania osobnej aplikacji zawierającej bibliotekę do uczenia maszynowego (Scipy), przestarzały wygląd UI
    \item \textbf{Model biznesowy:} Bezpłatna 
\end{itemize}

\newpage
\section{TrebEdit} 
\MyFigure[0.3\textwidth]{images/konkurencja/trebedit.png}{Aplikacja TrebEdit}{fig:trebedit}{Google Play}
\begin{itemize}
    
    \item \textbf{Docelowa platforma:} Android
    \item \textbf{Docelowa grupa:} Początkujący web-deweloperzy, osoby uczące się poprzez analizę cudzego kodu 
    \item \textbf{Główne przewagi konkurencyjne:} Możliwość sklonowania dowolnej strony WWW do edycji i eksperymentów, dedykowany pasek skrótów zawierający najpopularniejsze tagi i znaki specjalne, wbudowane mini-kursy i przypomnienia składni dla HTML, CSS, JS oraz PHP, zintegrowana konsola JavaScript i podgląd na żywo.
    \item \textbf{Słabe strony:} Brak wsparcia dla backendu, brak integracji z Git, aplikacja potrafi mieć problemy z asocjacją plików lub ich prawidłowym zapisem. 
    \item \textbf{Model biznesowy:} Bezpłatna z reklamami lub wersja Premium (oferująca m.in. sugestie kodu, sugestie ścieżki plików, emulację urządzenia, niestandardową czcionkę edytora, możliwość edycji paska narzędzi edytora, nielimitowaną ilość otwartych kart, brak reklam): jednorazowa opłata za 169,99 zł lub subskrypcja: miesięczna za 10,99 zł lub roczna za 99,99 zł. 
\end{itemize}
\section{Wnioski}
Z powyższej analizy wynika, że aplikacje naszej konkurencji są często skomplikowane, drogie i nierzadko wykorzystują uciążliwe systemy subskrypcji. Biorąc to pod uwagę, uważamy, że najlepszą decyzją we wczesnej fazie rozwoju projektu będzie skupienie się na głównym module aplikacji - narzędziach do tworzenia stron internetowych. Monetyzacja to również ważny temat - tu ustaliliśmy, że głównym źródłem finansowania będą dobrowolne dotacje od użytkowników. Przewidujemy także jednorazową opłatę za pobranie aplikacji ze sklepu App Store (iOS), jednak będzie to kwota wysoce przystępna.


\addcontentsline{toc}{chapter}{Projektowanie}
\chapter{Pomysł na Aplikacje}

Aplikacja mobilna „CodeMobile” to narzędzie stworzone z myślą o programistach frontendowych, którzy cenią sobie możliwość pracy z dowolnego miejsca na świecie. Nazwa aplikacji stanowi proste połączenie słów "Code" (kod) oraz "Mobile" (mobilny), co bezpośrednio oddaje jej główny cel: dostarczenie zaawansowanych funkcji znanych z desktopowych środowisk IDE w wygodnej, przenośnej formie.
\section{Logo}
\MyFigure[0.5\textwidth]{images/favicon.png}{Logo apilkacji}{fig:icon}{Opracowanie własne}
Ikona aplikacji \ref{fig:icon} nawiązuje kolorystyką do popularnego edytora kodu visual studio code, stworzonego przez Microslop przed erą AI. Ikona próbuje w ten sposób komunikować główny cel aplikacji którym jest bycie solidnym edytorem kodu na urządzenia mobilne.
\newline\newline
W odróżnieniu jednak do VSC aplikacja stara się być prosta a jednocześnie dostarczać potężne rozwiązania pomagające w tworzeniu stron internetowych w językach HTML, CSS oraz Javascript.

\section{Podstawowe funkcje}
Celem naszego projektu było stworzenie aplikacji z której sami chcielibyśmy korzystań. Jako studenci informatyki mamy już doświadczenie z programowaniem dlatego wiemy jakie funkcje faktycznie mogą nam się przydać.
\newline\newline
Mamy spore doświadczenie z programowaniem zarówno na komputerach tu możemy się pochwalić stroną internetową \url{https://iaml.uk/} oraz kontem github \url{https://github.com/thenameisluk} . W tym paroma osobistymi projektami takimi jak nasz makieta którą pisałem ręcznie bez użycia AI.
\newline\newline
Również makieta naszej aplikacji została napisana w czystym html'u, css'ie oraz javascript'cie również bez wykorzystania AI czyli dokładnie tej samej technologi w której programowanie aplikacja ma obsługiwać.
\newline\newline
Na urządzeniach mobilnych używałem między innymi Scriptable (\url{https://scriptable.app/}) oraz Codeapp (\url{https://apps.apple.com/au/app/code-app/id1512938504}) w którym stworzyłem między innymi taką prostą strone internetową \url{https://thenameisluk.github.io/Skill-Issue/} .
\newline\newline
Dlatego ze względu na nasze doświadczenie w tej dziedzinie uznaliśmy że posiadamy kompetencje żeby zaprojektować aplikacje do programowania z której faktyczni ludzie chcieli by korzystać. Bez reklam oraz płatności żeby pozwolić ludziom stworzyć coś ciekawego bez potrzeby płacenia lub korzystanie z słabszych aplikacji które są na rynku.
\newline\newline
Skompilowaliśmy więc listę funkcji które nasza aplikacja powinna posiadać żeby być użyteczna. Nie są to wszystkie funkcje które chcielibyśmy dodać ale dodanie wszystkiego co przyszło by nam do głowy sprawiło by że nasza aplikacja nie była by prosta w obsłudze wybraliśmy więc odpowiednią ilość funkcji żeby nie przesycić aplikacji a jednocześnie pozwolić użytkownikowi zrobić wszystko czego będzie potrzebować.
\subsection{Lsp}
Tak zwany language server provider \parencite{lsp} to protokół stworzony przez Microslop we współpracy z Red Hat oraz Codenvy. Protokół pozwala na integracje serwera językowego z dowolnym edytorem kodu wspierającym protokół, w ten sposób każdy nowy edytor kodu nie musi wynajdować koła na nowo dodając wsparcie dla pojedynczego języka.
\newline\newline
Serwer językowy może dostarczać klientowi takie funkcje jak wykrywanie błędów w kodzie, rodzaje tokenów w kodzie (znane jako syntax highlighting) oraz możliwe uzupełnienia kodu w pozycji kursora. Dzięki temu nasza aplikacja powinna być w stanie wyłapać błędy w kodzie bez potrzeby jego uruchamiania oraz umożliwić inteligentne uzupełniania kodu co jest naszą kolejną funkcją.
\subsection{Automatyczne uzupełnianie}
\MyFigure[0.5\textwidth]{images/konkurencja/scriptable.png}{Zrzut ekranu z aplikacji Sriptable}{fig:scriptable}{\url{https://www.youtube.com/watch?v=CjTzhPxfz8Q}}
Scriptable to jedna z najbardziej nie doczenianych aplikacji na system IOS, jest ona prosta (zakładając że użytkownik zna język Javascript) a jednocześnie pozwala zrobić naprawde niesamowite rzeczy wykorzystówjąc api systemu IOS takie jak tworzenie widgetów czy złożonych powiadomień.
\newline\newline
Jednak funkcją która najbardziej wpadła mi w oko w tej aplikacji jest możliwość automatycznego uzupełniania kodu w mobilnym stylu uzupełniania tekstu. Aplikacja ta zawsze pokazuje dostępne uzupełnienia w jednym miejscu na ekranie \ref{fig:scriptable}, tuż nad klawiaturą w odróżnieniu do textastic, androidIDE czy spck które uzupełnianie kodu oferują w miejscu kursora.
\newline\newline
Rozwiązanie to jest o wiele bardziej przyjemne i przyjazne dla użytkowników już przyzwyczajonych do korzystania z automatycznego uzupełniania w klawiaturze mobilne i wywołuje wrażenia aplikacji napisaną z myślą mobile first zabiast desktop first.
Naszym celem nie jest dostarczenie tradycyjnego edytora kodu na platforme mobilną dla nas najważniejsze jest żeby apilkacja była ergonomiczna i wygodna na platforme na którą ją wypuszczamy biorąc pod uwage ograniczenia jak i zalety plaformy.
\subsection{Git}
Git to najpopularniejszy system kontroli wersji na pisany przez Linus'a Torvalds'a twórce systemu Linux. W odróżnieniu do prostych użytkowników systemów Windows czy Mac OS, Linus wiedział że system kontroli wersji to jedno z najważniejszych narządzi podczas programowania dlatego ponieważ w momencie tworzenia pierwszych wersji jego systemu żadne narządzie na rynku nie odpowiadało jego wymaganiom Linus stworzył własne narzędzie Git.
\newline\newline
Od tego czasu git stał się jednym z najważniejszych narzędzi prawdziwych informatyków na całym świecie, platformy takie jak github, gitlab oraz gitea to głównie hosting'i git'ów ale również w pewnym sensie media społecznościowe dla deweloperów, pozwalające im dzielić się swoim kodem, rozwiązywać problemy jak i współpracować.
\newline\newline
W obecnym świecie git jest bardzo potężnym narzędziem i tworzenia projektów bez niego nie jest nie możliwe ułatwia nam monitorowanie progresu naszego kodu oraz zarządzanie nim. Choć nasza aplikacja nie będzie udostępniała zaawansowanych funkcji git'a (było by to nie możliwe jeśli chcemy zachować prostotę naszego interfejsu). Chcemy udostępnić naszym użytkownikom podstawowe funkcje git'a (to jest commit, push oraz fetch)
\subsection{Dokumentacja}
Dokumentacja to jedno z najważniejszych źródeł do których można mieć dostęp podczas programowania. Jest to dokument opisujący dostępne dla programisty funkcje oraz ich działanie. Choć słynny Rosyjski programista znany pod pseudonimem Tsoding z dostępem do prywatnej wersji beta do nie wypuszczonego jeszcze języka JAI uważa że czytanie kodu źródłowego języka jest o wiele lepszą techniką uczenia się języka lub szukanie funkcji których chce się użyć.
\newline\newline
Zmuszenie naszych użytkowników do czyta czytanie kodu źródłowego V8 lub Spidermonkey (silników Javascript napisanych w większości w C++) mogło by się skończyć ich utratą. Dlatego mimo rady tak istotnego programisty nasza aplikacja będzie zawierałą dokumentacje dostępnych funkcji i klas w języku Javascript (w ramach przeglądarki internetowej) oraz właściwości css.

\section{Funkcje których nie będzie}
\subsection{Terminal}
Choć mocna nad tym ubolewam terminal nie będzie częścią pierwszej wersji naszej apikacji, ponieważ nasza aplikacja celuje w prostotę a niektórzy nasi użytkownicy mogą być nie zaznajomieni z konceptem terminala nie chcieli byśmy ich zastraszyć tak skomplikowaną funkcją.
\newline\newline
Mimo że terminal mógł by pozwolić użytkownikom wykorzystać pełnie funkcji programu git było by to mniej intuicyjne niż używanie naszego interfejsu graficznego, nasza aplikacja skierowana jest w kierunku programowania stron internetowych a jedyną funkcją która mogła by wykorzystywać terminal przy tym zastosowaniu jest właśnie git.
% -*- TeX-master: "../main.tex" -*-
\chapter{Projekt Interfejsu}
Prace nad interfejsem rozpoczęliśmy od przygotowania wstępnego szkicu aplikacji \ref{fig:init-proj}, wykorzystując do tego narzędzie Discord Whiteboard.
\newline\newline
\MyFigure{images/initial.png}{Początkowy projekt}{fig:init-proj}{Discord Whiteboard}
Dzięki temu szkicowi dokładnie wiedzieliśmy, jakie ekrany będzie zawierać nasza aplikacja oraz jak będą wyglądały przejścia między nimi. Kluczową kwestią była dla nas spójność interfejsu. Kierowaliśmy się przy tym zasadą: za każdą konkretną czynność powinien odpowiadać wyłącznie jeden dedykowany element na ekranie.
\section{Pop-up}
Do zaprezentowania użytkownikowi predefiniowanych opcji wykorzystaliśmy okna typu pop-up. W naszej aplikacji służą one między innymi do wyboru metody importu danych (rysunek \ref{fig:import_screen}) czy określenia rodzaju konta (rysunek \ref{fig:settings_screen_acc}).
\MyFigure[0.3\textwidth]{images/elementy/popup.png}{Przykład Pop-up'u w naszej aplikacji.}{fig:popup}{\parencite{makieta}}
\section{Jumper}
Podczas projektowania interfejsu często wykorzystywaliśmy szuflady wysuwające się z dołu ekranu (nazywanych w naszej makiecie jumperami \ref{fig:jumper}). Jest to alternatywa dla znanych chociażby z systemu Android okien dialogowych, która - w naszej opinii - jest bardziej przyjazna dla użytkownika.
\MyFigure[0.3\textwidth]{images/elementy/jumper.png}{Przykład jumper'a w naszej aplikacji.}{fig:jumper}{\parencite{makieta}}
Ten element UI wysuwa się od spodu, delikatnie przyciemniając tło. Wykorzystywany jest przez nas głównie do pobierania dodatkowych danych od użytkownika (rysunek \ref{fig:new_proj_screen}) oraz wyświetlania informacji, którymi jest zainteresowany w danej chwili, takich jak na przykład dokumentacja (rysunek \ref{fig:code_doc_screen}).
\newline\newline
Ogromną zaletą tego rozwiązania jest intuicyjny sposób zamykania. Wystarczy, że użytkownik przesunie palcem w dół lub dotknie przyciemnionego obszaru poza panelem. Jest to dobry sposób na uniknięcie pomyłkowego kliknięcia. Co więcej, panel nie zasłania środkowej części ekranu, a jedynie jego dolny fragment. Dzięki temu użytkownik, czytając dokumentację, nie traci z oczu kodu, który właśnie pisał.
\newline\newline
Rozwiązanie to jest również wygodniejsze dla osób obsługujących aplikację jedną ręką, ponieważ pozwala na zamknięcie okna bez konieczności niewygodnego sięgania palcem do przycisku "anuluj". W efekcie nasza aplikacja jest nie tylko przystępniejsza dla szerszego grona odbiorców, ale również jej prezentacja jest o wiele nowocześniejsza niż konkurencyjne rozwiązania, takie jak na przykład Carnets.
\section{Slider}
Ostatnim istotnym elementem interfejsu naszej aplikacji, o działaniu zbliżonym do jumpera, jest tzw. slider \ref{fig:slider}. Elementy tego typu są często spotykane chociażby w aplikacjach na system Android. Slider funkcjonuje niemal identycznie co jumper, z tą różnicą, że wysuwa się z boku ekranu.
\MyFigure[0.3\textwidth]{images/elementy/slider.png}{Przykład slider'a w naszej aplikacji.}{fig:slider}{\parencite{makieta}}
Wykorzystujemy ten element jako substytut osobnego okna aplikacji. Pozwala on na szybki podgląd informacji o projekcie bez konieczności opuszczania aktualnego widoku, takich jak struktura plików czy błędy w kodzie.
\newline\newline
Dzięki połączeniu tych dwóch elementów (slider \ref{fig:slider} i jumper \ref{fig:jumper}), użytkownik programując cały czas pozostaje w jednym oknie edytora. Opuszcza je wyłącznie wtedy, gdy chce sprawdzić stan swojego projektu.


\addcontentsline{toc}{chapter}{Budowa}
% -*- TeX-master: "../main.tex" -*-
\chapter{Szczegółowy opis aplikacji}

Nasza aplikacja jest z założenia minimalna i elegancka dlatego taką właśnie ją zaprojektowaliśmy. Podczas wstępnego projektowania (rysunek \ref{fig:init-proj}) nie zajmowaliśmy się takimi szczegółami jak konkretna czcionka czy konkretny kolor interfejsu. Ale kiedy doszliśmy już do wstępnego projektu nadszedł czas żeby dopracować szczegóły zaczynając od czcionki.

\section{Dobór czcionki}

Jako główną czcionkę naszej aplikacji postanowiliśmy zastosować znaną otwarto źródłową czcionkę Intern (\url{https://github.com/rsms/inter}) jest to jedna z najlepszych prostych czcionek na rynku stworzona z pasją i delikatnością wykorzystywana między innymi przez takie projekty jak GNOME, Gitlab czy Unity (silnik do tworzenia gier).
\newline\newline
Natomiast jako czcionkę edytora tekstu domyślnie wykorzystujemy JetBrains Mono (\url{https://www.jetbrains.com/lp/mono/}) jedną z najpopularniejszych czcionek dla programistów. Czcionki mono albo monospaced są to specjalne czcionki często wykorzystywane do wyświetlania kodu, ponieważ charakteryzują się one tym że każdy znak w nich zawarty ma taką samą szerokość co pozwala edytorom jak i programistą wyrównywać do siebie tekst i tworzyć ładne ascii art'y bez potrzeby martwienia się szerokością konkretnego znaku. JetBrains Mono posiada również dodatkową cechę która czyni ją często wybieraną czcionką a mianowicie zamiana kombinacji znaków takich jak == lub != w jeden dłuższy znak po szczegóły polecam zapoznać się z stroną serwującą czcionkę.

\section{Ekran główny}

\MyFigure[0.5\textwidth]{images/ekrany/glowny.png}{Ekran główny aplikacji}{fig:main_screen}{\parencite{makieta}}
Pierwszym ekranem który zobaczą użytkownicy jest ekran główny (rysunek \ref{fig:main_screen}), jego zadaniem jest powitanie użytkownika w aplikacji oraz pozwolenie mu przejść do projektu nad którym chciał by pracować.
\newline\newline
W pierwszej sekcji ekranu znajdują się 2 guziki pozwalając użytkownikowi utworzyć nowy projekt lub dodać/zaimportować już istniejący. W następnej sekcji natomiast mamy listę wszystkich ostatnio otwartych projektów co pozwala użytkownikowi na szybki dostęp do projektów które już zaczął. Wybór któregoś z tych projektów przenosi użytkownika bezpośrednio do edytora kodu.
\newline\newline
Dodatkowo na ekranie głównym znajdują się 2 dodatkowe ikonu, jedna to logo Ko-Fi (\url{https://ko-fi.com/}) w prawym górnym rogu przenoszące użytkownika na stronę pozwalającą mu przekazać darowiznę na projekt. druga to ikona ustawień w prawym dolnym rogu to ikona ustawień która przenosi użytkownika do ekranu ustawień.

\MyFigure[0.5\textwidth]{images/ekrany/nowy_projekt.png}{Ekran dodawania nowego projektu}{fig:new_proj_screen}{\parencite{makieta}}
Po kliknięciu przycisku nowego projektu użytkownikowi pokazywany jest jumper (rysunek \ref{fig:new_proj_screen}) z możliwością wpisania nazwy projektu oraz jego utworzenia. Po podaniu nazwy i kliknięciu utwórz użytkownik zostaje przeniesiony do edytora kodu.

\MyFigure[0.5\textwidth]{images/ekrany/import.png}{Ekran popup pozwalający importować projekt}{fig:import_screen}{\parencite{makieta}}
Wybranie opcji importowania projektu zaprezentuje użytkownikowi wybór (rysunek \ref{fig:import_screen}) między dodaniem projektu z plików lokalnych lub wykorzystanie istniejącego repozytorium Git. 
\newline\newline
Wybranie opcji dodania z lokalnych plików powinno przenieść użytkownika do systemowej aplikacji pozwalającej na wybranie folderu. Po wybraniu odpowiedniego folderu użytkownik zostaje przeniesiony do edytora kodu. Wygląd tego meny zależy od wariantu oraz wersji systemu ze względu na to nie został on uwzględniony w naszej makiecie a wybranie tej opcji skutkuje przeniesieniem do edytora kodu.

\MyFigure[0.5\textwidth]{images/ekrany/import_git.png}{Ekran pozwalający importować projekt git}{fig:import_git_screen}{\parencite{makieta}}
Jeśli użytkownik zdecyduje się na importowanie istniejącego repozytorium git, zostanie mu wyświetlony jumper z możliwością podanie adresu repozytorium git oraz wybranie konta które ma zostać wykorzystane do pobrania danych z repozytorium. Konto git może zostać dodane do aplikacji z poziomu ustawień.

\section{Ustawienia}
\label{sec:settings}
\MyFigure[0.5\textwidth]{images/ekrany/ustawienia.png}{Ekran ustawień}{fig:settings_screen}{\parencite{makieta}}
Ekran ustawień (rysunek \ref{fig:settings_screen}) użytkownik może puścić klikajac strzałkę w górnym rogu ekranu. Ekran ten jest podzielony na trzy części, ustawienia edytora, ustawienia konta oraz ustatwienia dostępność. 
\newline\newline
W ustawiniech edytora użytkownik ma możliwość zmiany czcionki edytora oraz schematów kolorów opisanych w rozdziale \ref{sec:code_color}. Wybrany aktualnie kontrast jest zaznaczony niebieskim obramowaniem widocznym na rysunku.
\newline\newline
Kolejną opcją są ustawienia konta, pozwalające użytkownikowi na dodanie konta (więcej informacji na ten temat poniżej).
\newline\newline
Ostatnią kategorią jest kategoria dostępność, w której użytkownik może na chcile obecną uruchomić tylko jedną opcje dostępność którą jest tryb obsługi jedną ręką oraz wybrać rękę którą użytkownik będzie chciał obsługiwać aplikacja. Ta druga opcja pojawia się tylko w momęcie kiedy tryb obsługi jedną ręką jest włączony (pokazana na rysuneku \ref{fig:settings_screen}, ukryta na rysunku \ref{fig:settings_screen_acc}) żeby nie pokazywać użytkownikowi opcji która nic nie robi.

\MyFigure[0.5\textwidth]{images/ekrany/ustawienia_konto.png}{Ekran ustawień popup pozwalający na wybór konta}{fig:settings_screen_acc}{\parencite{makieta}}
\subsection{Konta serwisów git}
Wracając do opcji dodawanie konta użytkownik dostaje możliwość wyboru między trzema hosting'ami udostępniającymi serwisy git, mimo że trzy wydaje się sporą liczbą jest to nie wiele biorąc pod uwagę że niektóre projekty wolnego oprogramowania jaki i samodzielni deweloperzy często tworzą swój własnych server git do którego nie byli by się w stanie dostać do naszej aplikacji.
\newline\newline
Bardziej uniwersalnym rozwiązaniem było by wykorzystanie kluczy ssh do komunikacji z serwerami które są wspierane przez wszystkie nam znane serwisy git, jednak takie rozwiązanie było by o wiele bardziej skomplikowane i mniej intuicyjne dla zwykłych użytkowników którzy nawet nie wiedzą że git to nie jest github.
\newline\newline
Dodatkowym problemem było by faktyczne generowanie tych kluczy i zarządzanie nimi, klucze ssh zazwyczaj są generowane z poziomu komputera a następne zamieszczane na wybranym przez użytkownika serwerze. Jednak w przypadku aplikacji mobilnej zarządzanie tymi kluczami było by mniej intuicyjne i mogło by doprowadzić do flustracji użytkowników. Dlatego na chwile obecną zostajemy przy kontach do konkretnych usług ponieważ większa część użytkowników jest przyzwyczajona do ekranów logowania i wolała by się nie zajmować kluczami.

\section{Edytor kodu}
\label{sec:code}
\MyFigure[0.5\textwidth]{images/ekrany/kod.png}{Ekran edytora kodu}{fig:code_screen}{\parencite{makieta}}
Ekran edytora kodu jest najważniejszym ekranem w aplikacji, właśnie na tym ekrenie użytkownik będzie spędzał najwięcej czasu kożystając z aplikacji. Jest on minimalistyczny stara się pozostawić użytkownikowi jak najwięcej miejsca na faktyczny kod jednocześnie dając mu wystarczająco dużo informacji o jego stanie i tym gdzie się znajduje.
\newline\newline
W górnej cząści ekranu możemy zauważyć szary pasek zawierający informacje o pliku, to jest jego nazwą, typ, liczbę błędów oraz ostrzerzeń. W tym pask z lewej strony znajduje się charakterystyczna ikona trzech pasków pozwalająca przywołać slider z listą plików projektu oraz dodatkowymi opcjami, slider może być równeż wywołany poprzez przeciągnięcie od lewej strony krawiędzi.
\newline\newline
Poniżej znajduje się ekran podpowiedzi pokazujący użytkownikowi możliwe uzupełnienia aktyalniew spisywanego kodu, jest on inspirowany aplikacją scriptable (rysunek \ref{fig:scriptable}) a jego zawartość jest zarządzana przez lsp (więcej informacji w podrozdziale \ref{sub:lsp}).

\MyFigure[0.5\textwidth]{images/ekrany/kod_jb.png}{Ekran edytora kodu z włączonym trybem obsługi jedną ręką}{fig:code_jb_screen}{\parencite{makieta}}
Po włączeniu trybu obsługi jedną ręką, klawiatura ekranowa telefonu użytkownika powinna zostać przesunięta (rysunek \ref{fig:code_jb_screen}) bliżej krawędzi przy której znajduje się kciuk użytkownika, ułatwiając mu pisanie.

\MyFigure[0.5\textwidth]{images/ekrany/kod_doc.png}{Ekran dokumentacji edytora kodu}{fig:code_doc_screen}{\parencite{makieta}}
Będąc w edytorze kodu użytkownik ma możliwość przytrzymania (więcej informacji w podrozdziale \ref{sec:hold_sym}) wbudowanych funkcji wewnątrz kodu jak i podpowiedzi auto uzupełnienia żeby uzyskać dostęp do ich dokumentacji (rysunek \ref{fig:code_doc_screen}). Dokumentacja zawierać nazwę elemntu, jego typ oraz którki opis. Dokumentacja będzie dostępna dla użytkownika również w trybie offline ze względu na to że nasza aplikacja nie ma powodów żeby działać bez internetu a nasz zespół nie chce iść drogą aplikacji które odmawiają uruchomienia w momędzie kiedy urządzenie nie ma dostępu do sieci.

\MyFigure[0.5\textwidth]{images/ekrany/pliki.png}{Ekran plików edytora kodu}{fig:files_screen}{\parencite{makieta}}
Po wciśnięci ikonki trzech kresek lub przesunięciu od lewej strony użytkownikowi zostanie wyświetlony slider z animacją. Wewnątrz tego slidara znajdują się kluczowe elementy związane z kontrolą projektu w tym przycisk wyjścia z projektu w lewym górnym rogu, lista plików pozwalająca ukryć pliki wewnątrz katalogu klikając w nią oraz na samym dole kontrolki pozwalająde przejść do dodatkowych ekranów. Pierwsza kontrolka przenosi użytkownika do ekranu git'a, druga do ekranu podglądu projektu a trzecia do dokumentacji

\MyFigure[0.5\textwidth]{images/ekrany/bledy.png}{Ekran błędów edytora kodu}{fig:error_screen}{\parencite{makieta}}
Wciśnięcie ciemnego paska z ostrzerzeniam i błędami w kodie prezentuje użytkownika slider (rysunek \ref{fig:error_screen}) z szczegółowym opisem wszystkich błędów oraz linijką na której występują, błędy te są również podkreślane w edytorze kodu (rysunek \ref{fig:code_screen}).

\subsection{Dodawanie plików i folderów}
\MyFigure[0.5\textwidth]{images/ekrany/pliki_new.png}{Popup pozwalający na dodanie pliku lub folderu}{fig:files_new_screen}{\parencite{makieta}}
W menu z listą plików użytkownik ma możliwość przytrzymania na nazwe folderu projaktu albo inny folder plików celem otworzenie popup'u dodawania nowego pliku lub folderu (rysunke \ref{fig:files_new_screen}). Użytkownik ma opcje wybranie między dodaniem nowego pliku lub folderu po wybraniu użytkownikowi zostanie przedstawiony odpowiedni jumper (dla pliku rysunek \ref{fig:files_new_file_screen}, dla folderu rysunek \ref{fig:files_new_file_screen})

\MyFigure[0.5\textwidth]{images/ekrany/pliki_new_file.png}{Jumper pozwalający na podanie nazwy nowego pliku}{fig:files_new_file_screen}{\parencite{makieta}}

\MyFigure[0.5\textwidth]{images/ekrany/pliki_new_folder.png}{Jumper pozwalający na podanie nazwy nowego folderu}{fig:files_new_folder_screen}{\parencite{makieta}}


\section{Dodatkowe ekrany}

\MyFigure[0.5\textwidth]{images/ekrany/git.png}{Ekran git}{fig:git_screen}{\parencite{makieta}}
Ekran git'a (rysunek \ref{fig:git_screen}) pozwala użytkownikowi na dodanie commitów wraz z podaniem nazwy, wypchnięcie ich do zewnętrznego rezpozytorium git lub ściągnięcie z niego zmian. Użytkownik może również wybrać zmiany w których plikach zostaną zawarte w komicie, użycie git'a jest opcjonalne w naszej aplikacji ale dobrze integruje się z funkcją github'a pozwalającą hostować rezpozytoria git jako strony internetowe jak na przykład nasza makieta.

\MyFigure[0.5\textwidth]{images/ekrany/podglad.png}{Ekran podglądu}{fig:preview_screen}{\parencite{makieta}}
Ekran podglądu pozwala użytkownikowi podejrzeć wygląd strony nad którą pracuje, na dole ekranu znajduje się ikona powrotu oraz link wewnątrz strony na którym użytkownik obecnie się znajduje.

\MyFigure[0.5\textwidth]{images/ekrany/pomoc.png}{Ekran pomocy}{fig:help_screen}{\parencite{makieta}}
W ekranie dokumentacji lub pomocy (rysunek \ref{fig:help_screen}) znajduje się spis wszystkich klas dostępnych z dla użytkownika oraz ich metod. Wybranie którejść z metod kliknięciem otworzy jumper'a z dokumentacją (rysunek \ref{fig:help_doc_screen}) w takiej samej formie jak by się on pojawił na ekrani edytora co ma na zadanie za komunikować użytkownikowi że ta sama dokumentacja jest również dostępna z poziomu edytora i jeśli tego nie potrzbuje nie musi wchodzić do ekranu dokumentacji.
\MyFigure[0.5\textwidth]{images/ekrany/pomoc_doc.png}{Ekran dokumentacja pomocy}{fig:help_doc_screen}{\parencite{makieta}}

\chapter{Kontrast}
Kontrast to ważny temat dla naszej aplikacji, ponieważ cichlibyśmy żeby była ona dostępna dla jak największej ilości osób. Dlatego staramy się żeby tekst wyświetlany na ekranie był jak najlepiej widoczny.
\section{Ogólny kontrast}
Tła wewnątrz aplikacji maję 3 główne wartości (rysunek \ref{fig:bgpalet})

\begin{itemize}
    \item \textbf{Najciemniejszy} \#1e1e1e
    \item \textbf{Średnio ciemny} \#2f2f2f
    \item \textbf{Najjaśniejszy} \#585858
\end{itemize}

\MyFigure[0.5\textwidth]{images/kontrast/palataKolorówTła.png}{Paleta kolorów tła}{fig:bgpalet}{Ręcznie dobrane kolory}

Biały tekst na najjaśniejszym z tych kolorów posiada wartość 7.11 (rysunek \ref{fig:brightest}). Co spełnia zarówna wymagania AA jak i AAA.

\MyFigure[0.5\textwidth]{images/kontrast/brightest.png}{Biały napis na najjaśniejszym kolorze tła}{fig:brightest}{\url{https://colourcontrast.cc/?background=585858\&foreground=ffffff}}
Na ekranie głównym oraz w paru innych miejscach możemy róznież za obserwować tekst oraz guziki koloru niebieskiego (\#4b81ff). Ten kolor pojawia się jedynie na najciemniejszym tle a jego kontrast w tej sytuacji wynosi 4.67 (rysunek \ref{fig:bluetest}) co jest dopuszczalne w AA oraz AAA dla dużego tekstu.
\MyFigure[0.5\textwidth]{images/kontrast/bluetest.png}{Niebieski na najciemniejszym kolorze tła}{fig:bluetest}{\url{https://colourcontrast.cc/?background=585858\&foreground=ffffff}}

Ostatnią kombinacją kolorów w naszej aplikacji są niebierskie przyciuski z białymi napisami lub symbolami. Głównie w jumper'ach oraz na ekranie do zarządzania git'em.
\MyFigure[0.5\textwidth]{images/elementy/bluebutton.png}{Niebieski guzik na bialym}{fig:bluebutton}{\url{https://colourcontrast.cc/?background=325fc7\&foreground=ffffff}}

Ich kontrast wynosi 5.84 więc również jest git dla użych napisóœ zarówno w AAA jak i AA.
\MyFigure[0.5\textwidth]{images/kontrast/bluebutton.png}{Biały napis na niebieskim tle}{fig:bluebuttonk}{Ręcznie dobrane kolory}
\section{Kontrast Kodu}
Nasza aplikacja oferuje 4 schematy kolorów w edytorze kodu wszystkie prezentowane na tle \#1e1e1e. Schematy można wybierać z poziomu ustawień aplikacji. Każdy schemat ma swoje zastosowania i może być użyty przez różnych ludzi.
\subsection{Domyślny}
\MyFigure[0.5\textwidth]{images/schematy/default.png}{Schemat koloru domyślny}{fig:default}{\parencite{makieta}}
Podstawowy schemat kolorów inspirowany klasycznym motywem Dark+ w odróżnieniu do większości współczesnych motywów jest z tonowany i elegancki. Przypomina mi czasy kiedy zaczynałem moją przygodę z programowaniem.

\begin{table}[ht]
\centering
\caption{Schemat kolorów domyślny}
\label{tab:default}
\begin{tabular}{|>{\raggedright\arraybackslash}p{4cm}|c|c|}
\hline
\textbf{Element} & \textbf{Kolor Hex} & \textbf{Kontrast} \\
\hline
funkcja       & \#D8DD73 & 11.52:1 \\
operator      & \#BFBFAF & 8.96:1 \\
słowo kluczowe & \#E161D4 & 5.45:1 \\
zmienna       & \#8FC9EA & 9.3:1 \\
klasa         & \#48B988 & 6.8:1 \\
typ           & \#4795D9 & 5.21:1 \\
string        & \#C88B72 & 5.87:1 \\
stała         & \#9CDBF0 & 10.96:1 \\
liczba        & \#AAC295 & 8.62:1 \\
numer linii   & \#59B4DE & 7.15:1 \\
\hline
ostrzeżenie   & \#ECA25C & 7.85:1 \\
error         & \#ED9292 & 7.24:1 \\
\hline
\end{tabular}
\newline\newline
{\raggedright\fontsize{10pt}{12pt}\selectfont\textit{Źródło: Opracowanie własne}\par}
\end{table}

\subsection{Drugi}
\MyFigure[0.5\textwidth]{images/schematy/second.png}{Schemat koloru drugi}{fig:secong}{\parencite{makieta}}
Drugi motyw jest nieco bardziej edżi dla buntowników, emo lub ludzi bez smaku. Jakikolwiek balans jest wyrzucony za okno w imie jaskrawych kolorów bez żadnej klasy. Ludzie używający tego motywu nie wiedzieli by co to dobry motyw gdyby usiadł im na twarzy.

\begin{table}[ht]
\centering
\caption{Schemat kolorów drugi}
\label{tab:alt1}
\begin{tabular}{|>{\raggedright\arraybackslash}p{4cm}|c|c|}
\hline
\textbf{Element} & \textbf{Kolor Hex} & \textbf{Kontrast} \\
\hline
funkcja       & \#56B4E9 & 7.22:1 \\
operator      & \#CCCCCC & 10.38:1 \\
słowo kluczowe & \#E69F00 & 7.4:1 \\
zmienna       & \#56B4E9 & 7.22:1 \\
klasa         & \#009E73 & 4.87:1 \\
typ           & \#E69F00 & 7.4:1 \\
string        & \#F0E442 & 12.61:1 \\
stała         & \#F0E442 & 12.61:1 \\
liczba        & \#CC79A7 & 5.45:1 \\
numer linii   & \#59B4DE & 7.15:1 \\
\hline
ostrzeżenie   & \#ECA25C & 7.85:1 \\
error         & \#ED9292 & 7.24:1 \\
\hline
\end{tabular}
\newline\newline
{\raggedright\fontsize{10pt}{12pt}\selectfont\textit{Źródło: Opracowanie własne}\par}
\end{table}

\subsection{Daltonizm}
\MyFigure[0.5\textwidth]{images/schematy/dalton.png}{Schemat koloru dla daltonistów}{fig:dalton}{\parencite{makieta}}
Schemat kolorów oparty na palecie Bank Wong dla osób z daltonizmem.

\begin{table}[ht]
\centering
\caption{Schemat kolorów dla Daltonizm}
\label{tab:dalt}
\begin{tabular}{|>{\raggedright\arraybackslash}p{4cm}|c|c|}
\hline
\textbf{Element} & \textbf{Kolor Hex} & \textbf{Kontrast} \\
\hline
funkcja       & \#56B4E9 & 7.22:1 \\
operator      & \#CCCCCC & 10.38:1 \\
słowo kluczowe & \#E69F00 & 7.4:1 \\
zmienna       & \#80DDE0 & 10.58:1 \\
klasa         & \#00C792 & 7.6:1 \\
typ           & \#B8A2EC & 7.47:1 \\
string        & \#F0E442 & 12.6:1 \\
stała         & \#D897BB & 7.19:1 \\
liczba        & \#F4935A & 7.27:1 \\
numer linii   & \#AED8EF & 11:1 \\
\hline
ostrzeżenie   & \#FFD04D & 11.41:1 \\
error         & \#63D5D9 & 9.56:1 \\
\hline
\end{tabular}
\newline\newline
{\raggedright\fontsize{10pt}{12pt}\selectfont\textit{Źródło: Opracowanie własne}\par}
\end{table}

\subsection{Pusty}
\MyFigure[0.5\textwidth]{images/schematy/blank.png}{Schemat koloru pusty}{fig:blank}{\parencite{makieta}}
Cały tekst edytora jest biały dla maksymalnego kontrastu. Dla osób z ostrymi problemami z widzeniem.

\begin{table}[ht]
\centering
\caption{Schemat kolorów pusty}
\label{tab:bialy}
\begin{tabular}{|>{\raggedright\arraybackslash}p{4cm}|c|c|}
\hline
\textbf{Element} & \textbf{Kolor Hex} & \textbf{Kontrast} \\
\hline
wszystkie elementy & \#FFFFFF & 16.67:1 \\
\hline
\end{tabular}
\newline\newline
{\raggedright\fontsize{10pt}{12pt}\selectfont\textit{Źródło: Opracowanie własne}\par}
\end{table}
% -*- TeX-master: "../main.tex" -*-
\chapter{Dostępność}
Aby z naszego narzędzia mogła korzystać szeroka gama osób, wprowadziliśmy ustawienia dostępności, służące do dostosowania w zależności od potrzeb aplikacji. Zdajemy sobie sprawę z wielu wad genetycznych i różnych nieprzyjemności ograniczających możliwości użytkowników. Poniżej sporządziliśmy listę opcji, które sporządziliśmy dla naszego programu.

\section{Daltonizm}

Dużą częścią aplikacji jest kod, którego różne komponenty są kolorowane w zależności od ich zadania. Wiąże się to z możliwością problemu w odczycie dla użytkowników z wadami wzrokowymi, konkretniej z wadami uniemożliwiającymi odróżnianie niektórych barw. Na potrzeby tej grupy stworzyliśmy specjalny motyw na podstawie palety Wong (schemat \ref{tab:dalt}), która celuje w umożliwienie rodzajom daltonizmu takim protanopia, deuteranopia, tritanopia czy monochromatyzm czytelny odczyt.

\section{Tryb jednoręki}
Inną grupą osób z ograniczeniami, dla których chcieliśmy dostosować aplikację, są osoby, które nie są w stanie korzystać z klawiatury oburącz. Na ich potrzeby stworzyliśmy tryb jednoręki z wyborem na lewą lub prawą rękę. Po włączeniu klawiatura użytkownika przesuwana jest w wybraną stronę oraz lekko ścieśniana, co ułatwia wprowadzanie znaków przy pomocy wyłącznie jednej ręki.

% -*- TeX-master: "../main.tex" -*-
\chapter{Mikrointerakcje}
%To-do opisać mikrointerakcje

\MyFigure[0.5\textwidth]{images/elementy/redframe.png}{Czerwona ramka pojawiająca podczas próbu utworzenia projektu bez nazwy}{fig:redframe}{\parencite{makieta}}


\addcontentsline{toc}{chapter}{Zakończenie}
% -*- TeX-master: "../main.tex" -*-
\chapter{Makieta}
Kod źródłowy makiety \parencite{makieta} pod adresem \parencite{repo}. Została ona stworzona w czystym html'u, css'ie oraz javascript'cie bez używania żadnych frameworków ani AI w celu zachowania pełnej kontroli nad wyglądem makiety oraz animacjami.



\section{Symulacja przeciągania palcem}
Makieta pozwala na wykonywanie gestów przeciągnięcia myszką których obsługa została zaimplementowana w minimalnej funkcji.


\begin{lstlisting}[language=javascript, caption={Kod odpowiadający za gest przeciągania.}, label={code:przeciagniecie}]
function attachSwipeHandler(el=document.getElementById(""),fn=()=>{},right=true,corner=5,sens=3,updown=false){
    var lastP = 0;
    var eligible = false;
    el.addEventListener("mousedown" ,(/** @type {MouseEvent} **/ev)=>{
        if(lock)
            return;
        
        let p;
        if(!updown)
            p = ev.layerX*100/el.getBoundingClientRect().width;
        else
            p = ev.layerY*100/el.getBoundingClientRect().height;

        if(!right) p = 100-p;

        // console.log("md "+el.tagName+" "+Math.round(p));

        if(p<corner){
            // debugger
            console.log("eligible "+p+" < "+corner)
            eligible=true;
            lastP = p;
        }else eligible=false
        
    });
    el.addEventListener("mousemove" ,(/** @type {MouseEvent} **/ev)=>{
        if(lock)
            return;
        
        if(!updown)
            p = ev.layerX*100/el.getBoundingClientRect().width;
        else
            p = ev.layerY*100/el.getBoundingClientRect().height;

        if(!right) p = 100-p;
        
        // console.log("mm "+el.tagName+" "+Math.round(p));

        if(eligible){
            if(lastP+sens<=p){
                fn();
                eligible = false;
                console.log("triggered");
            }
        }
    });
    el.addEventListener("mouseup" ,(/** @type {MouseEvent} **/ev)=>{
        lock = false;
        // console.log("mu "+el.tagName+" ");
        eligible = false;
    });
}
\end{lstlisting}

Jak widać w kodzie \ref{code:przeciagniecie} funkcja działa wyłącznie dla myszki ponieważ w dla przeglądarki internetowej dotyk ekranu jest osobnym zbiorem eventów (to jest "touchstart" zamiast "mousedown") który wymaga wsparcia dla tak zwanego multitouch (po polsku wielo dotyku) co utrudnia prostą implementacje. Z tego powodu na laptopach z ekranami dotykowymi oraz na telefonach symulacja przeciągnięcia nie działa.
\newline\newline
Funkcje można wykorzystać do chowania slider'ów oraz jumper'ów zamiast klikania w przyciemnioną część ekranu. Dodatkowo pozwala również ona na przeciągnięcie od boku ekranu w edytorze kodu żeby wysunąć slider (rysunek \ref{fig:code_screen}) z listą plików w projekcie (rysunek \ref{fig:code_doc_screen}) bez potrzeby klikania przycisku hamburgera w górnym rogu ekranu.



\section{Symulacja wibracji urządzenia}
Makieta pozwala na symulacje wibracji telefonu po przez potrząśnięcie ekranu w momencie próby zatwierdzenia akcji bez wpisania informacji w polu tekstowym, funkcja ta jest wykorzystywana przy tworzeniu nowego projektu (rysunek \ref{fig:new_proj_screen}), dodawaniu repozytorium git (rysunek \ref{fig:import_git_screen}) oraz przy tworzeniu commit'a (rysunek \ref{fig:git_screen})

Kod który za to odpowiada jest jeszcze prostszy niż kot odpowiadający za przeciąganie.

\begin{lstlisting}[language=javascript, caption={Kod odpowiadający za wywoływania wibracji dodawanie czerwonej ramki do pustych pól.}, label={code:shake}]
function validateText(/**@type {Element}**/inp,fn=()=>{}){
    /**@type {string}**/
    let txt = inp.value
    
    if(txt==""){
        addTag(inp,"inp_wrong");
        addTag(app_screen,"shake")
        setTimeout(()=>{rmTag(app_screen,"shake")},400)
    }else{
        rmTag(inp,"inp_wrong");
        fn();
    }
}
\end{lstlisting}

\begin{lstlisting}[language=css, caption={Kod css odpowiadający za trzęsienie.}, label={code:shake-css}]
@keyframes shakeAnimation {
  from {transform: translate(-1vh,0);}
  to {  transform: translate(1vh,0);}
}

.shake{
  animation-name: shakeAnimation;
  animation-duration: 100ms;
  animation-iteration-count: infinite;
}
\end{lstlisting}

\section{Symulacja przytrzymania elementu}
\label{sec:hold_sym}
Makieta posiada również możliwość symulacji przytrzymania palcem elementu po przez kliknięcie i przytrzymanie klawisza myszki, jeśli klawisz zostanie podniesiony lub usunięty z guzika przed upływem odpowiedniej ilości czasu akcja zostanie anulowana.

Kod odpowiadający za tą funkcje załączam poniżej.

\begin{lstlisting}[language=javascript, caption={Kod odpowiadający obsługę przytrzymanie klawiszu myszki na elemencie.}, label={code:hold}]
function attachHoldHandler(el=document.getElementById(""),fn,wait=500){//500 milisekund
    var clicked = false;
    var id = 0;
    el.addEventListener("mousedown" ,(/** @type {MouseEvent} **/ev)=>{
        let remid = id;
        clicked = true;
        setTimeout(()=>{
            console.log(remid+" x "+id)
            if(clicked&&id==remid){
                fn()
            }
        },wait)
    })
    el.addEventListener("mouseleave",()=>{
        if(clicked){
            id++;
            clicked = false;
        }
    })
    el.addEventListener("mouseup",()=>{
        if(clicked){
            id++;
            clicked = false;
        }
    })
}
\end{lstlisting}

Jak widać w kodzie \ref{code:hold} przytrzymanie również można wywołać jedynie z poziomu myszki.



\chapter{Posdumowanie}

Projekt aplikacji "CodeMobile", pokazał nam ile możemy napotkać przeszkód decydując się wybór niekonwencjonalnych metod na projektowanie aplikacji, jednak że wyciągnęliśmy z tego zadania wiele ważnych wniosków jak na przykład to: jak ważne jest pozycjonowanie przycisków, dlaczego powinniśmy poświęcać więcej uwagi na kontrasty w naszej aplikacji i najważniejsza z tych wszystkich lekcja jak zaprojektować produkt, przy użyciu minimalnej ilości przycisków, tekstów i innych komponentów, aby nie przeciążyć użytkownika duża ilością informacji.
\section{Plany na przyszłość}
Jeśli mowa o przyszłości aplikacji, zastanawiamy się nad wypuszczeniem wersji na stystem IOS niestaty ta wersja będzie musiała byś płatna minimum 4 zł przy 100 użytkownikach ze względu na opłatę za konto developera w sklepie apple (około 400 zł rocznie).
\newline\newline
Dodatkowo jeśli środki na to pozwolą możemy dodać funkcjonalność terminalu dla git'a co umożliwiło by bardziej zaawansowanym użytkowniczką wykorzystać pełnie funkcji programu git.
\newline\newline
Innym pomysłem było by dodanie wsparcia dla języków kompilowanych jak C++ lub JAI (kiedy zostanie wypuszczony) która pozwalała by na pisanie prostych a jednocześnie mało zasobożernych gier mobilnych z użyciem bibliotek Raylib, SDL2 lub SFML.
\section{Zakończenie}
Na koniec najważniejsza lekcja o której będziemy pamiętać, prawdziwym skarbem nie była aplikacja lecz przyjaciele których spotkaliśmy po drodze i absolutna pasja która została wlana w kod naszej makiety <3.


% *************** Bibliografia ***************
% wymagania - Praca musi zawierać co najmniej 25 pozycji literaturowych, w tym 4 pozycje obcojęzyczne. Przypisy muszą być zgodne ze stylem APA. jako {nazwisko, rok, strona}
\newpage
\chapter*{Literatura} % Rozdział bez numeru
\addcontentsline{toc}{chapter}{Literatura} % Dodanie do spisu treści

\section*{Bibliografia}
\setcounter{mybibcounter}{0} 
\printbibliography[heading=none, notkeyword=netografia] % Wyświetla źródła nieoznaczone jako netografia

\section*{Netografia}
\setcounter{mybibcounter}{0} 
\printbibliography[heading=none, keyword=netografia] % Wyświetla źródła oznaczone jako netografia


% *************** Spisy pomocnicze ***************
\clearpage
\phantomsection
\addcontentsline{toc}{chapter}{Spis rysunków}
\listoffigures

\clearpage
\phantomsection
\addcontentsline{toc}{chapter}{Spis tabel}
\listoftables

% Jeśli chcesz dodać spis kodów (próbek kodu) 
\clearpage
\phantomsection
\addcontentsline{toc}{chapter}{Spis listingów}
\lstlistoflistings

\end{document}


